% =============================================================================
% Estudio económico de inferencia LLM local — prompt v4 (principle-driven)
% Compatible con el main_overleaf_llm.tex / memoria UPV-EARTH.
% Requiere en el preámbulo: \usepackage{booktabs, array, siunitx, graphicx}
% =============================================================================

\section{Estudio económico de la inferencia LLM local}
\label{sec:estudio-economico}

\subsection{Motivación}

La elección del modelo base (Anexo~0.1.1) se justificó en términos de
calidad de clasificación: \emph{flexible agreement}, \emph{rigorousness}
y \emph{positivity bias} medidos sobre el subset validado. Esa decisión
deja, sin embargo, una pregunta operativa abierta: \textbf{¿cuánto cuesta
realmente clasificar el corpus con cada modelo?}. En un despliegue real
del sistema UPV--EARTH ese coste no es despreciable, porque la
clasificación se ejecuta sobre los $\sim 31{,}600$ documentos del corpus
institucional ampliado y debe re--ejecutarse cuando cambia el prompt o
el catálogo de PBs. Este estudio cuantifica tres dimensiones del coste de
ejecutar el \textbf{prompt v4} (\emph{principle-driven} con tres
calibration cases, ver Sección~\ref{sec:prompt-v4}) en los tres modelos
open-weight de la comparativa M2:

\begin{enumerate}
  \item \textbf{Coste en tokens}: tokens de entrada y de salida por documento.
  \item \textbf{Coste en tiempo}: \emph{wall--clock} de cada llamada.
  \item \textbf{Coste energético}: potencia GPU integrada en el tiempo de
        inferencia, traducida a coste eléctrico y comparada con la tarifa
        equivalente si los mismos tokens se facturasen vía API en la nube.
\end{enumerate}

\subsection{Protocolo experimental}

\paragraph{Hardware.} Las medidas se tomaron en la máquina de inferencia
del grupo, cuyas especificaciones se detallan en la
Tabla~\ref{tab:estudio-economico-hw}. Todos los modelos se ejecutan
localmente mediante \texttt{ollama serve} con cuantización por defecto
(Q4\_K\_M); ni la CPU ni el almacenamiento están en el camino crítico
del experimento, pero se documentan por reproducibilidad.

\begin{table}[h]
\centering\small
\caption{Especificación del banco de pruebas. La GPU es el recurso
limitante; CPU, RAM y almacenamiento solo entran en juego para servir
el corpus desde disco y orquestar las llamadas a Ollama.}
\label{tab:estudio-economico-hw}
\begin{tabular}{ll}
\toprule
Componente & Especificación \\
\midrule
GPU            & NVIDIA RTX 4500 Ada Generation, 24\,GiB GDDR6, TDP 210\,W \\
               & (Compute Capability 8.9, driver 580.126.09) \\
CPU            & Intel Core i9--14900K, 24 cores / 32 threads, hasta 6.0\,GHz \\
RAM            & 128\,GiB DDR5 \\
Almacenamiento & SSD NVMe Samsung 990 Pro 1\,TB (sistema y modelos Ollama) \\
               & $2\times$ HDD Seagate 20\,TB (corpus y backups) \\
SO             & Ubuntu 24.04.3 LTS, kernel 6.17.0 \\
Runtime LLM    & Ollama \texttt{serve} (cuantización Q4\_K\_M en GGUF) \\
\bottomrule
\end{tabular}
\end{table}

La GPU es el único recurso saturable del experimento: el resto del
sistema (CPU, RAM, NVMe) opera muy por debajo de su capacidad mientras
Ollama procesa los abstracts, lo que permite atribuir toda la variación
observada al modelo y al prompt.

\paragraph{Modelos.} Se comparan los mismos tres modelos open-weight
evaluados en M2:

\begin{itemize}
  \item \texttt{llama3.1:8b} -- 8\,B parámetros, $\sim 4.7$\,GiB en disco.
  \item \texttt{qwen2.5:14b} -- 14\,B parámetros, $\sim 9$\,GiB en disco.
  \item \texttt{gemma4:26b} -- 25.8\,B parámetros, $\sim 19$\,GiB en disco.
\end{itemize}

\paragraph{Muestra.} Se eligen los primeros 20 documentos del subset
validado (excluyendo los 3 calibration cases del prompt v4). El mismo
abstract se envía a los tres modelos con \emph{idéntico} prompt v4,
\texttt{temperature}=0 y \texttt{top\_p}=0.9.

\paragraph{Instrumentación.} Para cada llamada se registra:

\begin{itemize}
  \item Los contadores oficiales que devuelve Ollama:
        \texttt{prompt\_eval\_count} (tokens de entrada),
        \texttt{eval\_count} (tokens de salida) y los \texttt{duration}
        en nanosegundos (\texttt{load}, \texttt{prompt\_eval},
        \texttt{eval}, \texttt{total}).
  \item Un \emph{sampler} concurrente que invoca \texttt{nvidia-smi} cada
        $250$\,ms durante la inferencia, extrayendo
        \texttt{power.draw}, \texttt{memory.used} y \texttt{utilization.gpu}.
        La energía consumida se obtiene por integración trapezoidal
        de la potencia instantánea:
        $E_{\text{Wh}} \approx \tfrac{1}{3600}\sum_i \tfrac{1}{2}(P_i+P_{i-1})\,\Delta t_i$.
  \item Antes de cada modelo se ejecuta una llamada \emph{warm-up} con
        \texttt{num\_predict=1} para descontar el coste de carga del
        modelo en VRAM, que no representa coste recurrente en producción.
\end{itemize}

El \emph{wall-clock} de la primera llamada queda contaminado por la carga
en VRAM, pero se neutraliza con el warm-up. El consumo eléctrico de
referencia se calcula con una tarifa de \SI{0.15}{\eur\per\kilo\watt\hour}.
El equivalente cloud se calcula multiplicando tokens medidos por los
precios públicos de GPT--4o (\$2.50/\$10.00 por M tokens
in/out), Claude Sonnet (\$3.00/\$15.00) y Claude Haiku
(\$0.80/\$4.00), con una conversión de \$1 = \EUR{0.92}.

\subsection{Resultados}

\begin{table}[h]
\centering\small
\caption{Coste de inferencia del prompt v4 por modelo. Tokens y tiempo
son medias por documento; energía y coste agregan los 20 documentos.
La columna \emph{Éxito} indica cuántos documentos completaron sin error
de \emph{timeout}.}
\label{tab:estudio-economico-modelos}
\begin{tabular}{lccccccc}
\toprule
Modelo & Éxito & tokens$_\text{in}$ & tokens$_\text{out}$ & wall (s) & $\bar{P}$ (W) & VRAM (GiB) & $E_\text{total}$ (Wh) \\
\midrule
\texttt{llama3.1:8b} & 20/20 & 2{,}391 & 107.2 & 2.78 & 112.8 & 10.5 & 1.66 \\
\texttt{qwen2.5:14b} & 20/20 & 2{,}423 & 105.3 & 4.08 & 139.0 & 16.2 & 3.04 \\
\texttt{gemma4:26b}  & 14/20$^\dagger$ & 2{,}503$^\ddagger$ & 149.5$^\ddagger$ & 1.78$^\dagger$ & 134.8 & 20.4 & 0.90$^\dagger$ \\
\bottomrule
\end{tabular}
\par\vspace{2pt}
\footnotesize
$^\dagger$ Excluye 6 documentos en los que la generación de gemma4:26b
agotó el \emph{timeout} de 600\,s.\quad
$^\ddagger$ El contador de tokens lo devolvió únicamente $1$ de las $14$
llamadas exitosas; el resto vino con \texttt{null} (Ollama no rellena
\texttt{prompt\_eval\_count} cuando hay \emph{cache hit} del prompt).
\end{table}

\begin{table}[h]
\centering\small
\caption{Coste por documento traducido a euros: electricidad local
(tarifa \SI{0.15}{\eur\per\kilo\watt\hour}) y precio equivalente si los
mismos tokens se facturaran en una API cloud (USD/M tok convertido a
EUR a \$1=\EUR{0.92}).}
\label{tab:estudio-economico-coste}
\begin{tabular}{lcccc}
\toprule
Modelo & € luz/doc & € GPT--4o/doc & € Claude Sonnet/doc & € Claude Haiku/doc \\
\midrule
\texttt{llama3.1:8b} & $1.25\!\times\!10^{-5}$ & $6.49\!\times\!10^{-3}$ & $8.08\!\times\!10^{-3}$ & $2.15\!\times\!10^{-3}$ \\
\texttt{qwen2.5:14b} & $2.28\!\times\!10^{-5}$ & $6.54\!\times\!10^{-3}$ & $8.14\!\times\!10^{-3}$ & $2.17\!\times\!10^{-3}$ \\
\texttt{gemma4:26b}  & $0.96\!\times\!10^{-5}$ & $6.79\!\times\!10^{-3}$ & $8.56\!\times\!10^{-3}$ & $2.28\!\times\!10^{-3}$ \\
\bottomrule
\end{tabular}
\end{table}

Los tres modelos consumen un orden de magnitud parecido en tokens de
entrada ($\sim 2{,}400$, dominado por las reglas de las 9 PBs que se
inyectan en el prompt v4) y entre $105$ y $150$ tokens de salida (un JSON
estructurado con razonamiento corto, PB primaria, PBs secundarias y
rechazadas). El coste energético local es, en todos los casos, despreciable
frente al coste equivalente cloud: \textbf{ejecutar el clasificador
localmente cuesta entre $260$ y $700$ veces menos electricidad que
facturar los mismos tokens a Claude Haiku}, el modelo cloud más barato
de la comparación.

\subsection{Lectura por modelo}

\paragraph{\texttt{llama3.1:8b}: el más eficiente, el menos riguroso.}
Es el más rápido (2.78\,s/doc), el de menor VRAM (10.5\,GiB, cabe en una
GPU de consumo) y el de menor potencia media (113\,W). Sin embargo, M2
ya documentó que tiene \emph{rigorousness} $0.00\,\%$: clasifica casi
cualquier abstract como PB y no aprende a abstenerse. El bajo coste
energético no compensa esa carencia operativa.

\paragraph{\texttt{qwen2.5:14b}: el compromiso defendible.}
Tarda ~$1.5\times$ más que llama y dobla su consumo energético, pero
exhibe la mejor combinación de \emph{flexible agreement} ($52\,\%$) con
\emph{rigorousness} apreciable ($38\,\%$) y \emph{positivity bias} $0.00\,\%$
(Anexo~0.1.1). El coste real por documento es ínfimo ($2.3\!\times\!10^{-5}$\,€
de electricidad) y la VRAM se mantiene cómodamente por debajo del límite
de la 4500 Ada. Es el único modelo de los tres que es a la vez
\emph{operacionalmente estable} y \emph{científicamente honesto}, lo que
justifica retrospectivamente la decisión de iterar el prompt v1$\to$v4
sobre qwen.

\paragraph{\texttt{gemma4:26b}: no es viable a 26\,B en esta GPU.}
Este es el hallazgo más relevante del estudio. La tabla
\ref{tab:estudio-economico-modelos} muestra una \emph{media}
engañosamente buena (1.78\,s/doc) porque solo agrega las llamadas que
completaron. La realidad cruda es:

\begin{itemize}
  \item \textbf{6 de 20 documentos ($30\,\%$) agotaron el timeout de
        600\,s}, es decir, gemma se quedó generando hasta que el cliente
        cortó la conexión. Cada uno de esos fallos consumió
        $\sim 25$\,Wh por sí solo, frente a los $0.06$\,Wh de una llamada
        normal.
  \item Las llamadas exitosas alternan entre tiempos muy bajos
        ($<1$\,s) que parecen \emph{cache hits} del prompt sin generación
        real, y respuestas razonables de $2$--$7$\,s. El contador
        \texttt{prompt\_eval\_count} viene a \texttt{null} en $13$ de las
        $14$ llamadas exitosas, lo que confirma que Ollama está sirviendo
        respuestas con el prompt cacheado y no re-evaluando.
  \item Pico de VRAM $20.4$\,GiB sobre $24$\,GiB disponibles: gemma vive
        en el límite del hardware, lo que explica la inestabilidad.
\end{itemize}

En términos de coste real --- incluyendo los timeouts --- gemma4:26b
consumiría aproximadamente $\mathbf{150}$\,\textbf{Wh por cada 20
documentos}, frente a $3$\,Wh de qwen, es decir, $\sim 50\times$ más
energía con $30\,\%$ de tasa de fallo. La afirmación de M2 sobre que
gemma era \emph{``demasiado lento para iterar sobre el corpus''} queda
ahora cuantificada con un argumento energético adicional: no solo es
lento, es inestable en esta GPU.

\subsection{Extrapolación al corpus completo}

Suponiendo coste lineal en el número de documentos y manteniendo el
prompt v4, clasificar los $31{,}634$ documentos del corpus institucional
ampliado con \texttt{qwen2.5:14b} costaría aproximadamente:

\begin{itemize}
  \item \textbf{Tiempo}: $31{,}634 \times 4.08\,\text{s} \approx 36$ horas
        de inferencia continua.
  \item \textbf{Energía}: $31{,}634 \times 0.152\,\text{Wh} \approx 4.8$\,kWh.
  \item \textbf{Coste eléctrico}: $\approx \EUR{0.72}$.
  \item \textbf{Coste equivalente cloud}: $\EUR{205}$ con GPT--4o,
        $\EUR{257}$ con Claude Sonnet, $\EUR{68}$ con Claude Haiku.
\end{itemize}

Es decir, el sistema completo se reclasifica end-to-end por menos de un
euro de electricidad, lo que hace operacionalmente trivial re-correr el
pipeline cada vez que el catálogo de PBs o el prompt sufran un cambio
significativo.

\subsection{Limitaciones del estudio}

\begin{enumerate}
  \item La muestra es de 20 documentos por modelo. Es suficiente para
        comparar órdenes de magnitud, pero las medias por documento
        tienen ruido del orden del $10$--$20\,\%$.
  \item \texttt{nvidia-smi} reporta consumo total de la GPU, no solo el
        atribuible al proceso Ollama. En esta máquina no había otra carga
        en GPU durante el experimento, por lo que la atribución es
        razonable, pero no es estrictamente \emph{per-process}.
  \item Los precios cloud son los públicos de mayo 2026 y no incluyen
        descuentos por uso prolongado, prompt caching del proveedor ni
        latencia de red.
  \item El coste de \emph{warm-up} (carga del modelo en VRAM, $\sim 5$--$15$\,s
        en gemma) no se imputa al documento; en un servicio que ya tiene
        el modelo residente es despreciable, pero no lo sería en
        invocaciones esporádicas tipo \emph{cold start}.
\end{enumerate}

Los CSV brutos del experimento (una fila por (modelo, documento) y
resumen agregado por modelo) se incluyen como material reproducible en
\path{nlp/llm/outputs/studies/estudio_economico_v4/}.
